\subsection{问题三：光伏预报、多时刻调整与偏差分段结算}

\subsubsection{问题分析}

问题三在问题二的基础上引入两级信息落差，二者分别触及模型的输入与结构。一是粒度落差：光伏预报只在每天 0:00、6:00、12:00、18:00 发布，且只给整点值，可决策与结算的粒度为 10 分钟，因此必须先把整点预报降到 10 分钟。二是时刻落差：可以按预报调整当天计划，但计划量与调整量之间要按“低于计划只退 0.5 倍价、高于计划需补 1.5 倍价”双向分段结算，因此必须明确在什么时刻、依据什么信息调整。本文据此先明确各决策时刻各自能看到什么，再建立“0:00 计划 + 三个时刻调整 + 事后结算”的三层模型，并保留紧急购电作为最终兜底。

本题很容易被忽略的一点是：调整机制并非没有代价。它减少了紧急购电，却引入了偏差违约金，二者孰大孰小只能由定量结果回答——这正是后文给出费用分项的原因。

\subsubsection{预报降尺度与信息集}

本文的做法是：把“预报 $k$ 小时”解释为发布时刻 $m$ 之后第 $k$ 个整点的功率点值；时段 $i$ 的左端点落在小时区间 $[h,h+1)$ 时取 $P(h)$ 与 $P(h+1)$ 的线性插值；发布时刻之后的第一小时无左端整点值，用 $P(m+1)$ 前向保持。把小时级预报细化到分钟级的插值方法及其对调度结果的影响在相关文献中已有系统讨论\cite{ref-almpantis-2024-minute-downscaling,ref-barbieri-2017-very-short-term-pv,ref-stenzel-2016-temporal-resolution}。

该做法有两个必须说明的性质。其一，\textbf{不强制块内能量守恒}：整点锚定线性插值得到的块内平均功率为 $(2.5A_{k-1}+3.5A_k)/6$，一般不等于 $A_k$，因此降尺度会系统性地改变能量。其二，\textbf{跨年项不消费}：$m=6,12,18$ 时分别有 19--24、13--24、7--24 小时的预报落在次日之后，合计 36 项被丢弃，本文不做截断也不做外推。这两点的影响在模型族对比中单独量化。

信息集按“模型假设”一节的非预期性约定逐层落实，此处只补充两点实现后果：附件 2 的实际光伏只进入事后结算层，不进入任何决策层；时段 $i$ 由唯一一个决策时刻支配，记为 $\nu(i)=6\lfloor(i-1)/36\rfloor$，无需引入离散选择变量，模型保持线性。

\subsubsection{三层模型的建立}

按时间顺序组织模型为三层，逐层求解、逐层提交，不允许用一次性联合优化让前两层“看到”实际光伏。

\textbf{计划层（0:00）}使用 0:00 预报的全天光伏轨迹 $\Pi_0[PV]$，求解

\begin{equation}
\min_{b,\,c,\,q^{dis},\,s,\,E}\ \sum_{i=1}^{144}p_ib_i
\quad\text{s.t.}\quad
b_i+\Pi_0[PV]_i\Delta t+q^{dis}_i=L_i\Delta t+c_i+s_i
\end{equation}

外加储能动态、容量与功率约束；储电量在日边界衔接上一日已提交的末态，全年终端自由。

\textbf{调整层（6:00、12:00、18:00）}以当天 0:00 的计划量 $b$ 为参数，只决定其支配时段的最终购电量 $q$，目标为该层剩余时段上的偏差惩罚：

\begin{equation}
\min\ \sum_{i\in R_m}\Bigl[\beta_{def}\,p_i\,\bigl(b_i-q_i\bigr)^{+}+\beta_{over}\,p_i\,\bigl(q_i-b_i\bigr)^{+}\Bigr],
\qquad \beta_{def}=0.5,\ \ \beta_{over}=1.5
\end{equation}

其中 $(x)^+=\max(x,0)$，实现时用两个非负辅助变量分段线性化，因系数为正故线性松弛无损；调整层初值取实际已实现的储电量，已提交时段不再修改。这类“滚动调整 + 双向偏差惩罚”的多阶段结构在日内再调度与偏差结算文献中被广泛采用\cite{ref-li-2016-two-stage-rhc,ref-bhattacharya-2018-multistage-storage,ref-elkazaz-2020-convex-mpc-rhc,ref-danti-2019-intraday-rescheduling,ref-matsumoto-2022-imbalance-design,ref-asgarpoor-1995-deviation-penalty}。

\textbf{实际结算层}在全部决策提交后事后计算，是附件 2 实际光伏的唯一入口。设剩余量 $r=q+q^{dis}+PV^{act}\Delta t-L\Delta t-c$，则

\begin{equation}
s'=\max(0,r),\qquad q^{em}=\max(0,-r)
\end{equation}

\begin{equation}
C_{\text{plan}}=\sum p b,\quad
C_{\text{adj}}=\sum\Bigl[\beta_{def}p(b-q)^{+}+\beta_{over}p(q-b)^{+}\Bigr],\quad
C_{\text{em}}=\sum \alpha_{em}p\,q^{em}
\end{equation}

其中 $b$ 是计划购电量、$q$ 是最终购电量（替代量的含义），二者不可混用。

\subsubsection{主结果与交付结果}

交付期总费用为 $14\,540\,616.335652$ 元，构成为：计划购电费 $12\,362\,860.209641$ 元、调整结算费 $510\,208.940910$ 元、紧急购电费 $1\,667\,547.185102$ 元（全期三项分别为 $13\,885\,669.475555$、$563\,098.132533$、$1\,730\,408.537141$ 元）。费用落在独立构造的区间 $[6\,963\,514.20,16\,407\,319.63]$ 元内，自下界起位于 80.23\% 处；交付期均价 0.697655 元/kWh，落在最小电价与 5 倍最大电价之间。交付期紧急购电量 391\,274.485914 kWh、弃光量 2\,087\,727.847547 kWh、最终购电量 $\sum q=20\,842\,142.600809$ kWh（注意该值不是计划购电量 $\sum b=20\,359\,871.873818$ kWh）。

表~\ref{tab:p3a} 给出四个指定日期的全天费用及其分项与紧急购电的合并区间数和合计电量。四个日期的 0:00 与 24:00 储电量均为 1200.000000 kWh，充放电量按 4 小时块分别列示、不做冲抵。同样提醒，题面表 4 中的 2025 年 3 月 1 日及其区间为格式示例。

\begin{table}[htbp]
\centering
\small
\caption{问题三：四个指定日期的全天购电量、购电费分项与紧急购电结构（日期为列、指标为行，单位随各行标注）}
\label{tab:p3a}
\begin{tabular}{lrrrr}
\toprule
指标 & 2025-03-20 & 2025-06-21 & 2025-09-23 & 2025-12-21 \\
\midrule
全天购电量 $\sum q$ / kWh & 70\,040.887993 & 33\,532.164963 & 68\,433.174470 & 94\,749.308713 \\
全天购电费 / 元 & 52\,814.209286 & 22\,848.629380 & 47\,936.179117 & 62\,365.864792 \\
\quad 其中计划购电费 / 元 & 41\,095.385464 & 18\,055.815827 & 42\,772.756895 & 59\,764.834892 \\
\quad 其中调整结算费 / 元 & 2\,373.895890 & 1\,088.445153 & 0.000000 & 1\,182.684937 \\
\quad 其中紧急购电费 / 元 & 9\,344.927933 & 3\,704.368400 & 5\,163.422222 & 1\,418.344963 \\
\midrule
紧急购电合并区间数 / 天 & 5 & 7 & 9 & 8 \\
紧急购电量合计 / kWh & 2062.489153 & 840.356133 & 1150.194481 & 352.514922 \\
\bottomrule
\end{tabular}
\end{table}

\subsubsection{机制解释}

\textbf{调整费全部来自超发方向。} 式 (14) 中欠发方向费用为 $0.0$ 元、超发方向为 563\,098.132533 元，即全部调整成本来自“最终购电量高于计划”的时段（上调 5810 个时段、下调仅 86 个时段）。原因是计划层只看 0:00 预报而结算层使用实际光伏：预报低估光伏时计划量偏高（对应下调），高估时偏低（对应上调），实测以后者为主。这一不对称是本文刻意保留的：计划层的目标只含计划费用，不引入对紧急购电的预期，以保证与题面“根据 0:00 预报制定计划”的语义一致。

\textbf{紧急购电非零的成因是预报误差，而非购电上限。} 决策层用附件 3 预报、结算层用附件 2 实际值，二者的差异必然在结算层造成缺口，由 5 倍价的紧急购电补足，交付期合计 391\,274.485914 kWh。本问不设购电功率上限，紧急购电并非“上限被顶到”所致。

\textbf{问题三比问题二更贵，这是一个必须说明的结论。} 两问交付期净负荷\textbf{逐位相同}（17\,939\,189.494967 kWh），而交付期费用分别为 $14\,540\,616.34$ 元与 $12\,233\,050.83$ 元，差额为

\begin{equation}
2\,307\,565.50\ \text{元}\quad(+18.863\%)
\end{equation}

构成为紧急购电费 1\,667\,547.185102 元加调整费 510\,208.940910 元，本质是“预报误差的代价”。本文据此明确指出：\emph{问题三并不比问题二更省}。模型之所以没有抵消这一代价，是因为计划层的目标 $\min\sum p\,b$ 只看 0:00 预报、对紧急购电没有任何预期，模型内部缺少对冲机制；这一事实同时构成了后文讨论“是否需要引入随机优化”的必要性证据。

\textbf{日边界储电量恒为下限，且这是结构性结果而非题面要求。} 全年 365 天的 $E_{d,144}$ 全部等于 1200 kWh，唯一取值数为 1（图~\ref{fig:p3soc}）。题面只在问题一要求“0:00 与 24:00 储电量相同”，本问没有该要求；该现象源于“计划层为单日线性规划 + 目标只含计划费用 + 终端自由”这一组合使跨日库存的边际价值为零，属结构性涌现结果，不应表述为模型缺陷，也不应表述为“题面要求每天放空”。需要注意的是，日边界归零并不等于储能闲置：时段级解序列显示每个交付日内的储电量都覆盖了整个运行区间 $[1200,10800]$ kWh（日内极差中位数 9600 kWh），即储能仍在满量程吞吐，只是每日末回到下限。此外，2025 年 2 月 1 日 0:00 的储电量在本问为 1200.00 kWh、在问题二为 7950.00 kWh，差异源于两问信息集不同（问题二用实际值、本问为预报驱动的三层滚动），二者不可直接比较。

\begin{figure}[htbp]
\centering
\includegraphics[width=\linewidth]{fig13_p3_soc.png}
\caption{问题三日内储电量行为：(a) 代表日的满量程吞吐与日末回落到下限；(b) 全年每日覆盖全运行区间、日末恒为下限（365/365 天）}
\label{fig:p3soc}
\end{figure}

\begin{figure}[htbp]
\centering
\begin{minipage}[t]{0.485\textwidth}
\centering
\includegraphics[width=\linewidth]{fig10_p3_forecast.png}
\caption{光伏预报的降尺度与“预报—实际”落差}
\label{fig:p3a}
\end{minipage}\hfill
\begin{minipage}[t]{0.485\textwidth}
\centering
\includegraphics[width=\linewidth]{fig11_p3_profile.png}
\caption{计划量、最终购电量与紧急购电量的日内剖面（虚线为三个调整时刻）}
\label{fig:p3b}
\end{minipage}
\end{figure}

\subsubsection{模型族对比与关键建模选择的影响}

12 组对照的结果表明（表~\ref{tab:p3ab}），主结果不依赖实现路线与求解器设置：逐日解耦的重组实现与主模型逐位同值，两组独立第二实现的链级相对差不超过 $2.17\times10^{-5}$（约 350 元）。表中同时列出若干无收益、甚至方向与直觉相反的对照，本文按事先设定纪律如实记录、不作反向解读，各项影响在对应小节的对照表中给出。

\begin{table}[htbp]
\centering
\small
\caption{问题三的模型族与建模选择对照（全期 365 天，单位：元）}
\label{tab:p3ab}
\begin{tabular}{p{6.6cm}>{\raggedleft\arraybackslash}p{3.6cm}>{\raggedleft\arraybackslash}p{4.0cm}}
\toprule
对照项 & 费用 & 相对主模型 \\
\midrule
主模型（三层顺序向前递推 + 闭式结算） & 16\,179\,176.145228 & --- \\
逐日解耦重组实现 & 16\,179\,176.145228 & $0.000000\%$ \\
独立第二实现 & 16\,178\,825.507754 & $-2.17\times10^{-3}\%$ \\
完全信息上界（单一全年 LP） & 13\,758\,182.573724 & $-14.960\%$ \\
仅 0:00 计划、无调整 & 16\,505\,833.062742 & $+2.019\%$ \\
决策加密至每 3 小时 & 16\,179\,575.041038 & $+0.0025\%$ \\
决策加密至每 1 小时 & 16\,179\,766.212646 & $+0.0036\%$ \\
降尺度改为分段常数（前向保持） & 18\,741\,130.202668 & $+15.834\%$ \\
日初复位为 6000 kWh（非物理对照） & 15\,360\,260.582255 & $-5.061\%$ \\
\bottomrule
\end{tabular}
\end{table}

表~\ref{tab:p3ab} 中有三项建模选择的影响需要强调。第一，\textbf{降尺度处理方式的影响大于调整机制本身的价值}：改为分段常数前向保持使费用上升 2\,561\,954.057440 元（$+15.83\%$），其中紧急购电项贡献占 94.7\%，机制是该做法在主导时段低估光伏 2.68\%、在 12:00 预报域低估 11.29\%，从而系统性地制造缺口。第二，\textbf{完全信息上界比主模型低 14.96\%}，这是信息价值的度量，且该上界与问题二的 accepted 结果逐位相同——它同时充当问题三与问题二之间的独立交叉基准，说明两问的取数与时间对齐完全一致。第三，\textbf{调整机制的价值}为 326\,656.917514 元（相对“仅 0:00 计划”），由“减少的 889\,755.05 kWh 紧急购电”与“新增的 563\,098.13 元违约金”共同决定。

\subsubsection{稳健性，以及“是否需要引入其他时刻的预报”}

扰动实验表明：六组参数方向全部不反转，其中往返效率影响最大（单侧最大 10.9715\%），初始储电量对交付期费用\textbf{完全没有影响}（六档差异均为 0.0，已由独立重解逐位复现）；噪声族内部离散度小于 1\%，但期望费用随误差幅度系统性上移（$\sigma=5\%$ 约 $+3.9\%$、$\sigma=10\%$ 约 $+11.2\%$），原因是噪声只扰动结算层实际值而决策层仍使用预报，缺口由单侧紧急购电与非对称双向结算承接。唯一未通过的检验标准是“跨噪声幅度档的族间均值差 7.29\% 超过 3\%”，同幅度档内仅 0.06\%；该检验标准把“族内离散度”与“跨档均值一致性”绑在一起，属检验设计问题，不构成“模型对噪声不稳定”的证据。据此本问分级为“条件稳定（需给出适用边界）”。购电上限的可行分界落在 $(4375,5000]$ kW，终端条件的经济价值为 3325.625002 元（$+0.0228713\%$）。

关于题面“请分析是否需要引入其他时刻的预报制定调整购电策略”，本文的回答限定在\emph{决策频率}层面，并给出三层证据。第一层是“有没有调整”：加入 6:00、12:00、18:00 三个调整时刻可节省 326\,656.92 元（2.02\%），说明调整机制本身有效。第二层是“在四个时刻之上继续加密”：加密到每 3 小时或每 1 小时，费用反而分别上升 398.90 元与 590.07 元，即\textbf{没有可测收益}，而计算时间从 10.5 s 增加到 69.9 s 与 194.8 s。第三层是“新增预报信息量”：附件 3 只在 0:00、6:00、12:00、18:00 发布预报，本文不构造并不存在的输入，因此上述结论\textbf{不覆盖}“若能获得更频繁发布的预报是否有价值”这一问题——那属于信息供给的变化，需要新的数据来源才能回答。综合而言，在现有预报发布条件下，题面给定的四个决策时刻已经足够，无需引入其他时刻。
